% Working manuscript, not the ICLR 2027 submission template.
% No empirical improvements are claimed before actual experiments finish.
\documentclass[11pt]{article}
\usepackage[margin=1in]{geometry}
\usepackage{amsmath,amssymb,booktabs,hyperref}
\title{Learning What to Retain During Memory Consolidation}
\author{Anonymous authors}
\date{Working draft: September 3, 2026}
\newcommand{\method}{RetentionMem}
\begin{document}
\maketitle
\begin{abstract}
Long-running language agents repeatedly consolidate past interactions to operate
within bounded context and computation. Consolidation can erase distinctions
needed for factual recall or action selection, causing tasks that previously
succeeded to fail. We study consolidation-induced regression and propose a small
policy trained on paired task outcomes before and after memory changes.
\method{} selects among retaining existing memory, accepting a summary, and
augmenting a summary with residual source information. The policy estimates
regression risk from frozen representations and retention features, while a
budgeted selection procedure preserves source spans with high predicted benefit
per added token. Our evaluation protocol holds the writer, retriever and reader
fixed to isolate this mechanism, then tests transfer across compressors and
readers. It covers conversational question answering, agent trajectories,
legitimate information updates, repeated consolidation and amortized cost.
Empirical results will be added after the preregistered experiments are complete.
\end{abstract}

\section{Introduction}
Memory compression must decide which distinctions an agent can afford to lose.
For example, a history may state that a delivery address changed from Seoul to
Busan, followed by a planned business trip to Daejeon. A summary containing all
three city names can still lose the current address. Preserving a previous answer
without accounting for the update can instead retain an obsolete fact. The desired
object is the ability to answer correctly under the latest observations, including
the temporal or conditional scope of a historical fact.

Structured compression and retrieval planning offer efficient access to long
histories \cite{simplemem}. Topic-based buffering and offline maintenance separate
memory construction from online inference \cite{lightmem}. In action settings,
raw experiences and dynamically constructed strategies serve different purposes
\cite{adamem}. These designs motivate a local consolidation interface; they do
not establish that a proposed rewrite preserves downstream capabilities.

Our proposed contributions are (i) an operational measurement of task regressions
induced by consolidation, (ii) a lightweight risk predictor and residual-selection
policy, and (iii) a matched evaluation of accuracy and total resource use. These
are hypotheses to test, not empirical conclusions of this working draft.

\section{Consolidation-induced regression}
Let $H_t$ be the source history available at time $t$, $M$ its active memory, and
$M'$ a proposed local replacement. For a frozen reader $A$ and a correctness
predicate $J$ evaluated against the current history, define
\begin{equation}
r(q;M,M',H_t)=J(A(M,q),q,H_t)\,[1-J(A(M',q),q,H_t)].
\end{equation}
We measure conditional regression among questions that are correct before the
change. Questions already answered incorrectly are not counted as consolidation
regressions. Correctness is semantic and task-specific, not equality of answer or
action strings. A legitimate observation update changes the reference state.
Historical questions retain their specified temporal scope.

During training, source-derived diagnostic questions target rare details,
numeric and temporal distinctions, exceptions, relations and action conditions.
Each diagnostic has exact source offsets. A separate judge validates its target
against the full available source prefix; ambiguous targets are excluded. The
paired reader sees the memory and questions but not reference answers. Benchmark
questions and future observations are excluded from this process.

\section{Method}
\paragraph{Storage and local candidates.}
\method{} maintains an immutable archive and an active store. New observations
enter both immediately. At a session boundary or a buffer threshold, it selects
at most eight related active units. A frozen writer proposes a summary targeting
a $4{:}1$ reduction. The writer sees only the active region: repeatedly rebuilding
from the archive would conceal losses accumulated through earlier consolidation.
A separate generator derives at most eight diagnostics from the source prefix.
The archive can support maintenance but is not searched during primary QA.

\paragraph{Risk prediction.}
Frozen embeddings of the question, reference answer, source spans and candidate
are combined with explicit numeric, temporal and conditional-retention features.
The current implementation projects frozen embeddings with a fixed seed and
uses 267 features with hidden widths 256 and 64 and a sigmoid output. Only this
small MLP is trained. Its targets are observed regressions for summaries and
partially repaired summaries. Online maintenance does not replay the paired
reader; those calls are confined to offline label collection.

\paragraph{Residual selection.}
Let $p_\theta(q,M')$ denote predicted risk, $\tau$ the validation-selected
threshold, and $B$ a replacement budget. If the summary meets both constraints,
the policy accepts it. Otherwise, for high-risk diagnostics it considers adding
their source spans, merging overlapping intervals. It greedily selects additions
that reduce $\sum_q\max(p_\theta(q,M')-\tau,0)$ most per added token, subject to
$B$. When the repaired candidate still violates risk or size constraints, it
retains the original active region. Retention may leave the requested replacement
budget unmet; this is recorded instead of silently deleting information. Query
budgets remain hard limits.

\paragraph{Calibration and provenance.}
Training and validation are split by source family before generating variants.
The implementation selects the largest validation threshold whose accepted
candidates satisfy an empirical regression-rate target, with a minimum accepted
count. It returns a retain-all threshold when no setting qualifies. This is not a
distribution-free preservation guarantee. Source references are recorded, but
reference recall is not equivalent to retaining the actual answer-bearing span.

\section{Evaluation protocol}
\begin{table}[h]
\centering
\begin{tabular}{lll}
\toprule
Setting & Primary outcome & Resampling unit \\
\midrule
LoCoMo categories 1--4 & Category-specific F1, judge & Conversation \\
LongMemEval cleaned S & Official-rubric judge & Independent history \\
ALFWorld unseen & Environment success, actions & Environment \\
\bottomrule
\end{tabular}
\caption{Planned evaluation. No result values have been inserted.}
\end{table}

Policy training uses synthetic sources and ALFWorld training experiences only.
The pilot contains 500 source regions; full collection targets 3,000 training
and 750 validation regions. Three policy seeds are evaluated. Official LoCoMo
contains 1,540 questions in the primary four categories. LongMemEval uses its
cleaned 500-question S release, with full session history and assistant turns.
Question dates are taken from the dataset rather than the execution date.

Mechanism controls share the candidate writer, hybrid retriever and reader:
raw retrieval, unconditional consolidation, retention, lexical selection,
coverage/preservation/faithfulness selection, learned selection, and learned
selection with residuals. Raw retrieval and retention are operationally identical
in the current implementation and are not counted as independent replications.
Full-system baselines require audited original execution paths. In particular,
an implementation of public verification criteria is not labeled as the trained
TrustMem model \cite{trustmem}.

For ALFWorld, the intervention transforms training trajectory values while
preserving AdaMEM's retrieval keys, HNSW index, ordering and action/strategy flow.
The official environment success flag supplies the task outcome. Changed action
strings do not establish regressions. Environment-specific prerequisite-violation
measurement remains an implementation item before the full action analysis.

We compare memory query budgets of 2K, 4K and 8K and request repeated maintenance
passes at 1, 2, 4 and 8. Requested passes are reported separately from actual
consolidation lineage depth. A compression target is not a realized compression
ratio; matched-size analysis must use measured active tokens or equal replacement
ceilings. Transfer from structured to topical compression and from Qwen3-8B to
Qwen2.5-7B-Instruct holds the learned policy fixed. This tests two readers, not
generalization across all model families.

\section{Cost and failure analysis}
We separately record archive, active and index storage, writer and diagnostic
calls, embedding computation, policy inference, reader calls and evaluation
judges. Offline label collection and policy training are included when reporting
total deployment cost. For $N$ subsequent queries,
\begin{equation}
C(N)=C_{\mathrm{train}}+C_{\mathrm{write}}+C_{\mathrm{maintain}}+N C_{\mathrm{read}}.
\end{equation}
Components must use a common declared unit for break-even calculations. Token
volumes, wall time and monetary cost are not interchangeable. The implementation
records unknown price or missing-usage fields instead of imputing zero cost.
The principal figure should show accuracy against cumulative cost over repeated
maintenance, with matched-size controls and uncertainty grouped by conversation
or environment.

\section{Related work and limitations}
Repeated memory updates can damage useful knowledge \cite{faulty}; observing
damage alone is therefore insufficient novelty. Tiered systems can recover
information through provenance and source access \cite{tiermem}; the primary
experiment excludes query-time archive lookup to isolate selective retention.
TrustMem learns consolidation with explicit verification criteria \cite{trustmem}
and is the nearest comparison. The required distinction is whether task-outcome
supervision improves prediction and whether small residuals add value under equal
budgets. Our system does not guarantee arbitrary future-query preservation.
Diagnostics may miss capabilities; source targets and judges may be wrong;
policies may overfit a compressor or a narrow synthetic distribution. These
limitations must be tested rather than addressed by broader system claims.

\section{Results status}
The software pipeline has been exercised with deterministic fixtures and a mock
model API. Actual Qwen labeling, benchmark QA, original-system baselines and
ALFWorld unseen execution are pending. This manuscript intentionally contains no
accuracy improvement, statistical significance or acceptance prediction.

\bibliographystyle{plain}
\bibliography{references}
\end{document}
